\subsection{REGULAR ELEMENTS OF A LIE GROUP}

DEFINITION 2. An element of G is said to be regular if it is regular for the adjoint representation of G.

In other words (Prop. 1), an element \(g \in \mathbf{G}\) is regular if, for all elements \(g'\) in a neighbourhood of \(g\) in \(\mathbf{G}\) , the dimension of the nilspace of \(\operatorname{Ad}(g') - 1\) is equal to the dimension of the nilspace of \(\operatorname{Ad}(g) - 1\) .

PROPOSITION 2. Let G' be a finite dimensional Lie group over k and f an open morphism from G to G'. The image under f of a regular element of G is a regular element of G'. If the kernel of f is contained in the centre of G, an element \(g \in G\) is regular if and only if \(f(g)\) is regular.

Indeed, let \(\mathfrak{g}'\) be the Lie algebra of \(G'\) and \(\mathfrak{h}\) the ideal in \(\mathfrak{g}\) given by the kernel of \(Tf|\mathfrak{g}\) . Let \(\rho\) be the linear representation of \(G\) on \(\mathfrak{h}\) defined by \(\rho(g) = \operatorname{Ad} g|\mathfrak{h}\) for all \(g \in G\) , and let \(\operatorname{Ad} \circ f\) be the linear representation of \(G\) on \(\mathfrak{g}'\) given by the composite of \(f\) with the adjoint representation of \(G'\) . These linear representations define an exact sequence of \(G\) -modules:

\[
0 \to \mathfrak {h} \to \mathfrak {g} \to \mathfrak {g} ^ {\prime} \to 0.
\]

By Lemma 2, \(r_{\mathrm{Ad}} = r_{\rho} + r_{\mathrm{Ad}\circ f}\) . Since \(r_{\mathrm{Ad}\circ f} = r_{\mathrm{Ad}}\circ f\) and since \(f\) is an open map, \(r_{\mathrm{Ad}\circ f}^{0} = r_{\mathrm{Ad}}^{0}\circ f\) . Consequently:

\[
r _ {\mathrm{Ad}} - r _ {\mathrm{Ad}} ^ {0} = r _ {\rho} - r _ {\rho} ^ {0} + \left(r _ {\mathrm{Ad}} - r _ {\mathrm{Ad}} ^ {0}\right) \circ f.
\]

Thus, if g is regular, \((r_{\mathrm{Ad}} - r_{\mathrm{Ad}}^{0})(f(g)) = 0\) , which means that \(f(g)\) is regular. If the kernel of f is contained in the centre of G,

\[
r _ {\rho} (g) = r _ {\rho} ^ {0} (g) = \dim \mathfrak {h}
\]

for all \(g \in G\) . Consequently, if \(f(g)\) is regular, \(r_{\mathrm{Ad}}(g) = r_{\mathrm{Ad}}^{0}(g)\) , in other words, g is regular.

PROPOSITION 3. Let \(G_{1}\) and \(G_{2}\) be two finite dimensional Lie groups over k. An element \((g_{1}, g_{2})\) of \(G_{1} \times G_{2}\) is regular if and only if \(g_{1}\) and \(g_{2}\) are regular elements of \(G_{1}\) and \(G_{2}\) , respectively.

The condition is necessary by Prop. 2. We show that it is sufficient. For all \(g = (g_{1}, g_{2}) \in \mathrm{G}_{1} \times \mathrm{G}_{2}\) , \(r_{\mathrm{Ad}}(g) = r_{\mathrm{Ad}}(g_{1}) + r_{\mathrm{Ad}}(g_{2})\) . In view of Lemma 1 (ii), it follows that \(r_{\mathrm{Ad}}^{0}(g) = r_{\mathrm{Ad}}^{0}(g_{1}) + r_{\mathrm{Ad}}^{0}(g_{2})\) . If \(g_{1}\) and \(g_{2}\) are regular, \(r_{\mathrm{Ad}}^{0}(g_{1}) = r_{\mathrm{Ad}}(g_{1})\) and \(r_{\mathrm{Ad}}^{0}(g_{2}) = r_{\mathrm{Ad}}(g_{2})\) , so \(r_{\mathrm{Ad}}^{0}(g) = r_{\mathrm{Ad}}(g)\) , which means that \(g\) is regular.

Lemma 3. Let \(a \in \mathbf{G}\) and let \(\mathfrak{m}\) be a complement of \(\mathfrak{g}^1(a)\) in \(\mathfrak{g}\) . Let \(\mathbf{U}\) be a neighbourhood of 0 in \(\mathfrak{g}\) and exp an exponential map from \(\mathbf{U}\) to \(\mathbf{G}\) . The map

\[
f: (x, y) \mapsto (\exp y) a (\exp x) (\exp y) ^ {- 1}
\]

from \((\mathfrak{g}^1 (a)\times \mathfrak{m})\cap \mathrm{U}\) to G is étale at \((0,0)\) .

The tangent linear maps at 0 of the maps \(x \mapsto a(\exp x)\) and \(y \mapsto (\exp y)a(\exp y)^{-1}\) are the maps \(x \mapsto ax\) and \(y \mapsto ya - ay = a(a^{-1}ya - y)\) from \(\mathfrak{g}\) to \(T_aG = ag\) (Chap. III, §3, no. 12, Prop. 46). Consequently, the tangent map of \(f\) at \((0,0)\) is the map \((x,y) \mapsto ax + a(a^{-1}ya - y) = a(x + a^{-1}ya - y)\) from \(\mathfrak{g}^1(a) \times \mathfrak{m}\) to \(ag\) . This map is injective. Indeed, if \(x \in \mathfrak{g}^1(a), y \in \mathfrak{m}\) and if \(x + a^{-1}ya - y = 0\) , then \((\mathrm{Ad}(a) - 1)y = \mathrm{Ad}(a)x \in \mathfrak{g}^1(a)\) since \(\mathrm{Ad}(a)\mathfrak{g}^1(a) \subset \mathfrak{g}^1(a)\) . This implies that \(y \in \mathfrak{g}^1(a)\) and consequently that \(y = 0\) . Since \(\mathrm{Ad}(a)\) is injective on \(\mathfrak{g}^1(a)\) , it follows that \(x = 0\) . Since \(\dim \mathfrak{g} = \dim \mathfrak{g}^1(a) + \dim \mathfrak{m}\) , this shows that \(f\) is étale at \((0,0)\) .

PROPOSITION 4. Let \(a \in G\) and \(H\) be a Lie subgroup germ of \(G\) with Lie algebra \(\mathfrak{g}^1(a)\) . The map \((b, c) \mapsto cabc^{-1}\) from \(H \times G\) to \(G\) is a submersion at \((e, e)\) .

Indeed, let \(\mathfrak{m}\) be a complement of \(\mathfrak{g}^1(a)\) in \(\mathfrak{g}\) and exp an exponential map of \(G\) defined on an open neighbourhood \(U\) of \(0\) in \(\mathfrak{g}\) . We can choose \(U\) so that \(\exp(U \cap \mathfrak{g}^1(a)) \subset H\) . The map \(f: (x, y) \mapsto (\exp x, \exp y)\) is an analytic map on a neighbourhood of \((0, 0)\) in \(\mathfrak{g}^1(a) \times \mathfrak{m}\) with values in \(H \times G\) . By Lemma 3, the composite of \(f\) with the map \(\varphi: (b, c) \mapsto cabc^{-1}\) is étale at \((0, 0)\) . It follows that \(\varphi\) is a submersion at \(f(0, 0) = (e, e)\) .

PROPOSITION 5. Let \(a \in G\) and let \(W\) be a neighbourhood of \(e\) in \(G\) . There exists a neighbourhood \(V\) of \(a\) with the following property: for all \(a' \in V\) , there exists an element \(g \in W\) such that \(\mathfrak{g}^1(a') \subset \mathrm{Ad}(g)\mathfrak{g}^1(a)\) .

Put \(\mathfrak{g}^1 = \mathfrak{g}^1(a)\) and let \(\mathfrak{g} = \mathfrak{g}^1 + \mathfrak{g}^+\) be the Fitting decomposition of \(\operatorname{Ad}(a) - 1\) (§1, no. 1). Let H be a Lie subgroup germ of G with Lie algebra \(\mathfrak{g}^1\) . For all \(h \in \mathrm{H}\) , \(\operatorname{Ad}(h)\mathfrak{g}^1 \subset \mathfrak{g}^1\) . Since \([\mathfrak{g}^1, \mathfrak{g}^+] \subset \mathfrak{g}^+\) , there exists a neighbourhood U of \(e\) in H such that \(\operatorname{Ad}(h)\mathfrak{g}^+ \subset \mathfrak{g}^+\) for all \(h \in \mathrm{H}\) . Since the restriction of \(\operatorname{Ad}(a)-1\) to \(g^{+}\) is bijective, U can be chosen so that the restriction of \(\operatorname{Ad}(ah)-1\) to \(g^{+}\) is bijective for all \(h\in U\) . Then \(\mathfrak{g}^{1}(ah)\subset\mathfrak{g}^{1}(a)=\mathfrak{g}^{1}\) for all \(h\in U\) . By Proposition 4, \(\operatorname{Int}(\mathrm{W})(a\mathrm{U})\) is a neighbourhood of a in G. If \(a'\in\operatorname{Int}(\mathrm{W})(a\mathrm{U})\) , then \(a'=g(ah)g^{-1}\) with \(g\in W\) and \(h\in U\) ; it follows that \(\mathfrak{g}^{1}(a')=\operatorname{Ad}(g)\mathfrak{g}^{1}(ah)\subset\operatorname{Ad}(g)\mathfrak{g}^{1}(a)\) .

COROLLARY. Let \(G^{*}\) be an open subgroup of G. If \(a \in G\) is regular, there exists a neighbourhood V of a such that, for all \(a' \in V\) , \(\mathfrak{g}^{1}(a')\) is conjugate to \(\mathfrak{g}^{1}(a)\) under \(\mathrm{Ad}(\mathrm{G}^{*})\) .

